\documentclass[11pt]{exam}
\usepackage{foerderkurs}

\begin{document}

\minilektion{2.1}{Rechenreihenfolge und Klammern}
  {Terme in der richtigen Reihenfolge berechnen und erkennen, wann Klammern
   nötig sind.}

\begin{lehrperson}[Ablauf]
\textbf{Dauer:} ca.\ 12 Minuten Instruktion, danach Übungen in der Übungszeit,
Exit-Frage am Ende der Lektion.

\begin{tabular}{@{}ll@{}}
  3 min & Einstieg: Abstimmung zu beiden Aufgaben, kurze Diskussion \\
  4 min & Merkkasten mit den vier Stufen und den zwei Präzisierungen \\
  3 min & Beispiele 1–4 \\
  2 min & Bruchstrich, andere Schreibweisen, Auftrag für die Übungen
\end{tabular}

\textbf{Verknüpfung:} Standortbestimmung 2.4, 2.5, 9.3; Übungen mit den Tags
2B und 2C.

\textbf{Voraussetzung:} Mini-Lektion 0.1 (Gleichungsketten korrekt
aufschreiben).
\end{lehrperson}

% ---------------------------------------------------------------------------

\begin{einstieg}
\begin{enumerate}
  \item Diese Aufgabe wird im Internet regelmässig heftig diskutiert.
        Was ist das Ergebnis: $1$ oder $16$?
        \[ 8 \geteilt 2 \cdot (2 + 2) \]
  \item Was ergibt $18 - 6 \geteilt 2 + 1$? \quad
        \wahl{$7$}\quad \wahl{$14$}\quad \wahl{$16$}
\end{enumerate}
\end{einstieg}

\begin{lehrperson}[Einstieg]
Zuerst abstimmen lassen (Hände oder Whiteboards), dann nach Begründungen
fragen. Die Diskussion zu Aufgabe 1 auf etwa zwei Minuten begrenzen.

\textbf{Aufgabe 1:} Nach den Regeln ergibt sich $16$, weil Punktrechnungen
gleichrangig sind und von links nach rechts ausgeführt werden:
$8 \geteilt 2 = 4$, dann $4 \cdot 4 = 16$. Wer $1$ erhält, rechnet zuerst
$2 \cdot (2 + 2)$. Die eigentliche Botschaft: Der Ausdruck ist
\emph{schlecht geschrieben}. Wer etwas eindeutig mitteilen will, setzt
Klammern oder verwendet einen Bruchstrich.

\textbf{Aufgabe 2:} richtig ist $16$. Die $7$ entsteht beim Rechnen strikt von
links nach rechts, die $14$ bei «plus vor minus» ($18 - 3 - 1$ statt
$18 - 3 + 1$, also $18 - (3 + 1)$).
\end{lehrperson}

% ---------------------------------------------------------------------------

\begin{merke}
Gerechnet wird in dieser Reihenfolge:
\begin{enumerate}
  \item \textbf{Klammern}
  \item \textbf{Potenzen}
  \item \textbf{Punktrechnungen} ($\cdot$ und $\geteilt$)
  \item \textbf{Strichrechnungen} ($+$ und $-$)
\end{enumerate}
Zwei Präzisierungen, die oft vergessen werden:
\begin{itemize}
  \item Innerhalb einer Stufe sind die Operationen \textbf{gleichrangig} und
        werden \textbf{von links nach rechts} ausgeführt. «Mal vor durch» und
        «plus vor minus» gibt es nicht.
  \item Eine Potenz gilt nur für das, was \textbf{direkt davor} steht:
        $2 \cdot 3^2 = 2 \cdot 9 = 18$ \quad und \quad $-3^2 = -9$,
        aber \quad $(-3)^2 = 9$.
\end{itemize}
\end{merke}

% ---------------------------------------------------------------------------
\abschnitt{Beispiele}

\begin{beispiel}[1: Punkt vor Strich, dann von links nach rechts]
\[ 18 - 6 \geteilt 2 + 1 = 18 - 3 + 1 = 15 + 1 = 16 \]
\falsch\quad $18 - 3 + 1 = 18 - 4 = 14$ \quad («plus vor minus»)
\end{beispiel}

\begin{beispiel}[2: Mal und durch sind gleichrangig]
\[ 20 \geteilt 4 \cdot 5 = 5 \cdot 5 = 25 \]
\falsch\quad $20 \geteilt 4 \cdot 5 = 20 \geteilt 20 = 1$ \quad
(«mal vor durch»)
\end{beispiel}

\begin{beispiel}[3: Potenzen]
\begin{align*}
  2 \cdot 3^2    &= 2 \cdot 9 = 18 & (2 \cdot 3)^2 &= 6^2 = 36 \\
  -3^2           &= -9             & (-3)^2        &= (-3) \cdot (-3) = 9
\end{align*}
Bei $-3^2$ wird zuerst quadriert, dann kommt das Minus dazu. Soll die negative
Zahl quadriert werden, braucht es Klammern.
\end{beispiel}

\begin{beispiel}[4: Der Bruchstrich wirkt wie eine Klammer]
\[ \frac{9 + 15}{2} = \frac{24}{2} = 12
   \qquad\text{aber}\qquad
   9 + \frac{15}{2} = 9 + 7.5 = 16.5 \]
Schreibt man den Bruch auf einer Zeile, muss man die Klammer selbst setzen:
$\frac{9 + 15}{2} = (9 + 15) \geteilt 2$.
\end{beispiel}

\begin{lehrperson}[Beispiele]
Die beiden «falsch»-Zeilen in Beispiel 1 und 2 sind die häufigsten Fehler in
der Standortbestimmung (Aufgabe 2.4\,c und e). Sie entstehen oft aus der
Merkregel «Punkt vor Strich», die als «mal vor durch» und «plus vor minus»
missverstanden wird.

Bei Beispiel 3 hilft der Hinweis: $-3^2$ liest man «minus drei hoch zwei», also
«minus (drei hoch zwei)». Dieser Fehler führt später beim Einsetzen negativer
Zahlen zu falschen Ergebnissen (Mini-Lektion 5.4).
\end{lehrperson}

% ---------------------------------------------------------------------------

\begin{schreibweisen}
\begin{itemize}
  \item In englischsprachigen Videos heissen die Regeln \textbf{PEMDAS} oder
        \textbf{BODMAS}. Die Merkwörter ordnen die Buchstaben unterschiedlich
        an, aber die Regeln sind dieselben wie oben.
  \item Am Computer und in KI-Antworten schreibt man auf einer Zeile:
        \texttt{*} für mal, \texttt{/} für durch, \texttt{\^} für hoch.
        Dann ist $\frac{9+15}{2}$ \texttt{(9+15)/2}, aber \texttt{9+15/2}
        ist $9 + \frac{15}{2}$.
  \item Taschenrechner und Programme halten sich an dieselben Regeln, aber nur
        an das, was \emph{tatsächlich eingetippt} wird. \texttt{-3\^{}2} ergibt
        $-9$.
  \item Vorsicht: Einige Taschenrechner und Programme geben einer
        Multiplikation \emph{ohne} Malpunkt, wie in $2(2 + 2)$, Vorrang. Und
        in Tabellenkalkulationen wie Excel ergibt \texttt{=-3\^{}2} sogar $9$.
        Darum gilt: \textbf{Im Zweifel Klammern setzen.}
\end{itemize}
\end{schreibweisen}

% ---------------------------------------------------------------------------
\abschnitt{Übungen}

\begin{questions}

\question Berechnen Sie. Schreiben Sie die Zwischenschritte als korrekte
Gleichungskette auf.
\begin{multicols}{3}
\begin{parts}
  \part $7 + 3 \cdot 4$
  \part $(7 + 3) \cdot 4$
  \part $36 \geteilt 6 \geteilt 2$
  \part $36 \geteilt (6 \geteilt 2)$
  \part $24 - 8 + 4$
  \part $24 - (8 + 4)$
\end{parts}
\end{multicols}
\begin{solution}
  a) $19$ \quad b) $40$ \quad c) $3$ \quad d) $12$ \quad e) $20$ \quad f) $12$
\end{solution}
\raum{3.5cm}

\question Berechnen Sie.
\begin{multicols}{3}
\begin{parts}
  \part $3 \cdot 2^2$
  \part $(3 \cdot 2)^2$
  \part $-2^2$
  \part $(-2)^2$
  \part $-2^3$
  \part $10 - 2^3$
\end{parts}
\end{multicols}
\begin{solution}
  a) $12$ \quad b) $36$ \quad c) $-4$ \quad d) $4$ \quad e) $-8$ \quad f) $2$
\end{solution}
\raum{3cm}

\question Setzen Sie Klammern so, dass die Gleichung stimmt.
\begin{multicols}{2}
\begin{parts}
  \part $2 + 3 \cdot 4 = 20$
  \part $12 \geteilt 4 - 1 = 4$
  \part $5 + 3^2 = 64$
  \part $18 - 6 \geteilt 2 + 1 = 4$
\end{parts}
\end{multicols}
\begin{solution}
  a) $(2 + 3) \cdot 4 = 20$ \quad
  b) $12 \geteilt (4 - 1) = 4$ \quad
  c) $(5 + 3)^2 = 64$ \quad
  d) $(18 - 6) \geteilt (2 + 1) = 4$
\end{solution}
\raum{3cm}

\question
\begin{parts}
  \part Berechnen Sie: \quad $\dfrac{12 + 8}{4}$ \qquad und \qquad
        $12 + \dfrac{8}{4}$
  \part Schreiben Sie auf einer Zeile, so wie man es am Computer eintippen
        würde:
        \[ \frac{a + b}{2} \qquad\qquad \frac{3}{x + 1} \qquad\qquad
           \frac{x^2 - 1}{4} \]
\end{parts}
\begin{solution}
  a) $5$ und $14$ \quad
  b) \texttt{(a+b)/2}, \ \texttt{3/(x+1)}, \ \texttt{(x\^{}2-1)/4}
\end{solution}
\raum{3cm}

\question Lea rechnet:
\[ 5 + 2 \cdot 3^2 = 5 + 6^2 = 5 + 36 = 41 \]
Was hat sie falsch gemacht? Rechnen Sie richtig.
\begin{solution}
  Lea hat zuerst $2 \cdot 3$ gerechnet und dann quadriert. Die Potenz gilt nur
  für die $3$: \ $5 + 2 \cdot 3^2 = 5 + 2 \cdot 9 = 5 + 18 = 23$.
\end{solution}
\raum{2.5cm}

\end{questions}

\begin{lehrperson}[Übungen]
Übung 3 ist für viele die interessanteste, weil sie umgekehrt denken müssen:
Wo braucht es Klammern, damit das gewünschte Ergebnis entsteht? Teil d) ist
schwieriger, weil zwei Klammerpaare nötig sind. Übung 4\,b) bereitet das
Einsetzen in Formeln vor und ist für Schülerinnen und Schüler wichtig, die
viel mit KI-Programmen oder Tabellenkalkulationen arbeiten.
\end{lehrperson}

% ---------------------------------------------------------------------------

\begin{ausblick}
Die Rechenreihenfolge brauchen Sie bei jedem Term, beim Einsetzen von Zahlen
in Formeln (z.\,B.\ in der Physik) und bei jeder Eingabe in den Taschenrechner.
Viele Fehler in späteren Prüfungen entstehen nicht beim schwierigen Teil einer
Aufgabe, sondern bei einem Vorzeichen oder einer vergessenen Klammer.
\end{ausblick}

\begin{exitfrage}
Berechnen Sie: \quad $2 + 3 \cdot (4 - 1)^2$
\ifprintanswers\par\medskip\textbf{Lösung:}
$2 + 3 \cdot (4 - 1)^2 = 2 + 3 \cdot 3^2 = 2 + 3 \cdot 9 = 2 + 27 = 29$
\else\par\vspace{2cm}\fi
\end{exitfrage}

\end{document}
